\documentclass[fleqn,usenatbib]{mnras}

\usepackage[T1]{fontenc}
\usepackage{ae,aecompl}
\usepackage{newtxtext,newtxmath}
\usepackage{graphicx}
\usepackage{amsmath}
\usepackage{booktabs}
\usepackage{multirow}
\usepackage{xcolor}
\usepackage{hyperref}
\usepackage{lineno}
\linenumbers

\newcommand{\Vobs}{V_{\mathrm{obs}}}
\newcommand{\Vbar}{V_{\mathrm{bar}}}
\newcommand{\Vtot}{V_{\mathrm{tot}}}
\newcommand{\Vextra}{V_{\mathrm{extra}}}
\newcommand{\Leff}{L_{\mathrm{eff}}}
\newcommand{\Lzero}{L_0}

\title[Inner-region scatter response in SPARC]{Inner-Region Scatter Response to Bounded Perturbations: A Comparative Analysis of 74 SPARC Galaxies}

\author[V. M. F. Figueiredo]{
	Vitor M. F. Figueiredo$^{1}$\thanks{E-mail: vitor.figueiredo.research@protonmail.com}
	\\
	$^{1}$Independent Researcher, Nisa, Portugal
}

\date{Accepted XXX. Received YYY; in original form ZZZ}
\pubyear{2026}

\begin{document}
	\maketitle
	
\begin{abstract}
	We present a comparative analysis of inner-region scatter sensitivity 
	across three galaxy rotation-curve implementations using a common sample 
	of 74 galaxies from the SPARC database. The comparison is performed on a 
	harmonised basis, adopting the same inner-region definition 
	($R < 0.5\,R_{\max}$), the same operational scatter metric 
	(log$_{10}$-RMS residuals in dex), and bounded perturbation tests around 
	published or archived baseline solutions. The three cases considered are 
	Navarro--Frenk--White (NFW) halo models, the adopted radial acceleration 
	relation (RAR) implementation, and a fixed-parameter non-local kernel 
	implementation (H2).
	
	A clear ranked pattern emerges. The H2 implementation shows a 
	near-neutral response, with median 
	$|\Delta\sigma|_{\mathrm{H2}} = 0.000478$~dex on the full 74-galaxy 
	common sample. The adopted RAR implementation shows intermediate 
	sensitivity, with median $\max |\Delta\sigma|_{\mathrm{RAR}} = 0.0321$~dex 
	and no statistically significant monotonic correlation with the chosen 
	baryonic-dominance proxy ($\rho = +0.0485$, $p = 0.681$, $N = 74$). The 
	NFW baseline shows the largest response, with median 
	$\max |\Delta\sigma|_{\mathrm{NFW}} = 0.0548$~dex and a strong 
	anti-correlation with baryonic dominance 
	($\rho = -0.899$, $p = 1.8 \times 10^{-27}$, $N = 74$). These results 
	show that, under a common inner-region diagnostic, bounded perturbations 
	produce substantially different scatter responses across the tested 
	implementations, supporting the use of scatter sensitivity as a 
	comparative diagnostic for galaxy rotation-curve models. The results are 
	presented as a comparative diagnostic of implementation sensitivity 
	under bounded perturbations, rather than a fundamental test of the 
	underlying physical theories.
\end{abstract}
	
	\begin{keywords}
		galaxies: kinematics and dynamics -- galaxies: haloes -- dark matter -- gravitation -- methods: data analysis -- methods: numerical -- catalogues
	\end{keywords}
	
\section{Introduction}
\label{sec:intro}

The modelling of galaxy rotation curves remains a central tool 
for probing mass distribution and gravitational dynamics in disc 
galaxies. The observed discrepancy between baryonic mass 
distributions and measured circular velocities has traditionally 
been addressed through two broad classes of frameworks: dark-matter 
halo models within the standard cosmological picture, commonly 
represented by profiles such as Navarro--Frenk--White (NFW; 
\citealt{Navarro1997}), and modified-gravity or modified-inertia 
approaches, most notably Modified Newtonian Dynamics (MOND; 
\citealt{Milgrom1983}; for a review see \citealt{Famaey2012}) 
and its empirical phenomenology in the radial acceleration 
relation (RAR; \citealt{McGaugh2016, Lelli2017}).

Although these frameworks can often provide comparably good global 
descriptions of observed rotation curves, the inner regions of 
galaxies provide a particularly useful regime for comparative 
testing. In these high-acceleration regions, the baryonic 
contribution is frequently large, and bounded perturbations to 
any additional model component are expected to have more limited 
leverage on the total circular velocity. The diversity of 
inner rotation curve shapes across galaxy populations 
\citep{deBlok2010, Oman2015} further motivates examining model 
behaviour in this regime, making the inner rotation curve a 
natural setting for comparing not only model adequacy, but also 
model sensitivity.

A related issue is the extent to which bounded parameter variations 
can alter residual scatter in the inner regions. Existing comparative 
studies have primarily focused on best-fitting performance across 
different frameworks (e.g. \citealt{Lelli2016, Li2018, Li2020}), whereas 
less attention has been given to a harmonised comparison of bounded 
perturbation sensitivity under a common diagnostic. Such a comparison 
is useful because it helps separate structural response from the 
broader flexibility associated with unconstrained refitting.

In this work, we present a comparative inner-region scatter 
diagnostic applied to a common set of 74 galaxies from the SPARC 
database \citep{Lelli2016}. The comparison uses the same inner-region 
definition, the same operational scatter metric (log$_{10}$-RMS 
residuals in dex), and the same bounded-perturbation approach across 
three implementations: NFW halo models based on the catalogue of 
\citet{Li2020}, bounded sensitivity tests within the adopted RAR 
implementation, and a fixed-parameter non-local kernel implementation (H2).

We find a clear ranked pattern of inner-region scatter sensitivity 
across the tested implementations. The NFW case shows the largest 
sensitivity under bounded perturbations, with the response strongly 
modulated by baryonic dominance. The adopted RAR implementation 
shows intermediate sensitivity, while the H2 implementation exhibits 
near-neutral behaviour in the same diagnostic.

This paper is organised as follows. Section \ref{sec:methods} 
defines the operational diagnostic and its bounded-perturbation 
logic. Section \ref{sec:data} describes the SPARC data and the 
common comparison sample. Sections \ref{sec:nfw}, \ref{sec:rar}, 
and \ref{sec:h2} present the NFW, RAR, and H2 analyses, 
respectively. Section \ref{sec:comparative} presents the three-way 
comparison, Section \ref{sec:discussion} discusses the structural 
interpretation and limitations of the result, and Section 
\ref{sec:conclusions} summarises our conclusions.

\section{Methods: operational definitions}
\label{sec:methods}

Our aim is to compare the inner-region scatter response of three model
implementations under bounded perturbations using a common diagnostic.
To make this comparison operationally consistent, we adopt the same
inner-region definition, the same scatter metric, and the same
bounded-perturbation philosophy across all three test cases.

\subsection{Inner-region definition}
\label{subsec:inner_region}

For each galaxy, we define the inner region as
\begin{equation}
	R < 0.5\,R_{\max},
\end{equation}
where $R_{\max}$ is the largest radius sampled by the observed rotation
curve in the SPARC data set. This definition is used uniformly across
the NFW, RAR, and H2 analyses.

This choice is operational rather than model-dependent. It provides a
common radial domain that is straightforward to reproduce across all
implementations, while still focusing on the region in which the
baryonic contribution is often comparatively strong. Galaxies with
fewer than three valid inner-region data points are excluded from the
common comparison sample, as such cases do not provide a statistically
stable estimate of the adopted scatter metric.

\subsection{Inner-region scatter metric}
\label{subsec:scatter_metric}

For a given model realisation, we quantify the inner-region scatter
using the root-mean-square of the logarithmic velocity residuals,
\begin{equation}
	\sigma_{\mathrm{inner}}
	=
	\left[
	\frac{1}{N_{\mathrm{in}}}
	\sum_{i=1}^{N_{\mathrm{in}}}
	\left(
	\log_{10} V_{\mathrm{model},i}
	-
	\log_{10} V_{\mathrm{obs},i}
	\right)^2
	\right]^{1/2},
\end{equation}
where the sum runs over the $N_{\mathrm{in}}$ observed points satisfying
$R_i < 0.5\,R_{\max}$. The metric is therefore expressed in dex.

We adopt this form for two reasons. First, it is directly comparable
across all three model implementations used here. Secondly, the
logarithmic residual form measures relative rather than absolute
velocity offsets, which is preferable for a heterogeneous galaxy sample
spanning a broad dynamic range in circular velocity amplitude.

For each bounded perturbation, we evaluate the corresponding change in
inner-region scatter as
\begin{equation}
	\Delta \sigma
	=
	\sigma_{\mathrm{inner,\,perturbed}}
	-
	\sigma_{\mathrm{inner,\,baseline}}.
\end{equation}
For the NFW and RAR analyses, where multiple bounded perturbations are
applied around each galaxy's baseline model, we report the maximum 
absolute change, $\max |\Delta \sigma|$, as a compact summary of 
inner-region sensitivity. For the H2 implementation, the bounded 
adaptive test yields a single direct scatter response under the 
adopted perturbation protocol.

\subsection{Bounded perturbation framework}
\label{subsec:bounded_framework}

The comparative question addressed in this work is not which framework
achieves the globally best fit under unrestricted optimisation. Instead,
we ask how strongly the adopted inner-region scatter metric responds
when the non-baryonic or non-Newtonian part of the model is perturbed
within bounded ranges around an established baseline.

For a broad class of rotation-curve constructions, the model circular
velocity can be written schematically as
\begin{equation}
	V_{\mathrm{tot}}^2
	=
	V_{\mathrm{bar}}^2
	+
	V_{\mathrm{extra}}^2,
	\label{eq:quadrature_form}
\end{equation}
where $V_{\mathrm{bar}}$ denotes the baryonic contribution and
$V_{\mathrm{extra}}$ denotes an additional model-dependent component.
In baryon-dominated inner regions, where $V_{\mathrm{extra}} \ll
V_{\mathrm{bar}}$, bounded variations in the extra component are
expected to have reduced leverage on the total circular velocity. A
first-order sensitivity estimate makes this explicit:
\begin{equation}
	\delta V_{\mathrm{tot}}
	\approx
	\frac{V_{\mathrm{extra}}}{V_{\mathrm{tot}}}\,
	\delta V_{\mathrm{extra}},
	\label{eq:delta_vtot}
\end{equation}
so the response of $V_{\mathrm{tot}}$ to perturbations in the extra
component is suppressed as the ratio
$V_{\mathrm{extra}}/V_{\mathrm{tot}}$ becomes small.

Equation~(\ref{eq:delta_vtot}) does not by itself imply identical
behaviour across different model classes. Rather, it motivates the
working expectation that bounded inner-region scatter sensitivity should
depend on how strongly the additional component contributes within the
adopted radial domain. The present study tests whether this expectation 
holds across three distinct implementations when evaluated under a 
common diagnostic.

\subsection{Perturbation normalisation}
\label{subsec:perturbation_normalisation}

To keep the comparison operationally controlled, bounded perturbations
are applied around each implementation's baseline solution using fixed,
predefined ranges rather than unrestricted refits. For the NFW case,
halo parameters are perturbed by $\pm10$ and $\pm20$ per cent around
their catalogue values. For the adopted RAR implementation, the
perturbations are applied to the quantities specified in the comparison
catalogue and are interpreted strictly as bounded sensitivity tests
within that implementation. For H2, the perturbation protocol is fixed
by the bounded adaptive procedure defined in the archived analysis.
While the specific perturbation protocols differ in
their operational details, they are not intended to be dynamically
equivalent across implementations. Different model architectures do
not admit a unique, model-independent measure of perturbation
amplitude. The objective is therefore not to impose identical
perturbation effort, but to evaluate how the same inner-region scatter
metric responds under bounded, finite variations defined within the
operational limits of each implementation’s baseline solution.

These bounded perturbations are not intended to represent exhaustive
re-optimisation of each framework. Their purpose is narrower: to test
how strongly the same inner-region scatter metric responds to
controlled, finite variations in each implementation under otherwise
matched observational conditions.

\section{Data and sample selection}
\label{sec:data}

\subsection{SPARC rotation-curve data}
\label{subsec:sparc_data}

Our analysis is based on the SPARC database \citep{Lelli2016}, which
provides rotation curves and baryonic mass models for a large sample of
nearby disc galaxies. For each system, the database includes the
observed circular velocity profile together with the baryonic
decomposition into gaseous, stellar disc, and bulge components where
available. These data provide the common observational basis for all
three model implementations considered in this work.

We use the SPARC rotation-curve sampling directly, rather than
interpolating onto a new radial grid, so that all comparisons are
performed on the native observational radii of each galaxy. This
ensures that the inner-region scatter diagnostic is always evaluated on
the same observed points for a given system.

\subsection{Comparison catalogues}
\label{subsec:comparison_catalogues}

The comparative analysis combines the SPARC rotation-curve data with
two published external catalogues and one archived internal
implementation.

For the dark-matter halo baseline, we use the NFW fits reported by
\citet{Li2020}. The relevant catalogue quantities are the halo
parameters $V_{200}$, $c_{200}$, and related derived quantities under
the tabulated convention adopted by that work. These parameters define
the baseline NFW model around which bounded perturbations are applied.

For the modified-gravity comparison, we use the RAR catalogue reported
by \citet{Li2020}, which tabulates the best-fitting quantities for the
adopted RAR implementation on a galaxy-by-galaxy basis. In this case,
the catalogue provides the baseline implementation from which the
bounded sensitivity tests are constructed.

The third comparison case is the archived H2 fixed-parameter non-local
kernel implementation, which was previously evaluated on an 80-galaxy
SPARC fleet using the same operational inner-region scatter metric.
Only the harmonised inner-region diagnostic output of that analysis is
used here.

\subsection{Common comparison sample}
\label{subsec:common_sample}

The starting point for the present work is the 80-galaxy H2 fleet. We
then restrict the comparative analysis to galaxies for which all three
implementations can be evaluated under the common inner-region
diagnostic defined in Section~\ref{sec:methods}.

A galaxy is excluded from the final comparison if it has fewer than
three valid inner-region data points under the adopted cut
$R < 0.5\,R_{\max}$. This criterion is applied uniformly and identically 
across all three analyses in order to maintain direct comparability of 
the scatter metric. Six galaxies fail this requirement and are removed 
from the common set:
F567-2, NGC4389, NGC6789, UGC00634, UGC05999, and UGC09992.
The final 74-galaxy sample spans $\Vobs/\Vbar = 0.027$ to $2.187$, covering the full range from extreme baryon domination to dark-matter domination.

This quality-control criterion ($n_{\mathrm{inner}} \ge 3$) ensures that
the log$_{10}$-RMS scatter metric is statistically well-defined for every
galaxy in the common comparison set, preventing sparse inner-region
sampling from biasing the sensitivity diagnostic.

This common comparison set is used throughout the main comparative
analysis. Because the same six systems are excluded in all three
implementations, the final H2, NFW, and RAR results are evaluated on an
identical galaxy basis.

\subsection{Section ordering and comparative logic}
\label{subsec:section_ordering}

The three model implementations are presented in the order NFW, RAR,
and H2. This ordering is deliberate. The NFW case provides the most
familiar halo-based baseline within standard galaxy mass modelling and
therefore serves as the natural first comparison. The adopted RAR
implementation then provides a modified-gravity benchmark under the
same operational diagnostic. The H2 analysis is presented last so that
it can be interpreted as one tested implementation within the broader
comparative framework, rather than as the starting point of the paper.

This ordering is intended to keep the emphasis on the galaxy-model
comparison itself. The central question is not how the H2 framework
behaves in isolation, but how the inner-region scatter response differs
across three bounded model implementations when evaluated on the same
galaxies and under the same diagnostic conditions.
\section{NFW halo perturbation analysis}
\label{sec:nfw}

\subsection{NFW implementation}
\label{subsec:nfw_implementation}

For the dark-matter halo comparison, we use the NFW catalogue provided by
\citet{Li2020}. The baseline halo parameters are taken directly from the
tabulated best-fitting values, with the NFW implementation defined by the
catalogue quantities $V_{200}$, $c_{200}$, and the associated derived scale
radius $r_{\mathrm{s}}$. The adopted convention corresponds to a halo
overdensity of 200 times the critical density, using the cosmological
assumptions specified in the source catalogue.

For each galaxy, the halo contribution is evaluated on the native SPARC
rotation-curve radii described in Section~\ref{sec:data}. The total model
velocity is then constructed as
\begin{equation}
	V_{\mathrm{tot}}^2(R)
	=
	V_{\mathrm{bar}}^2(R)
	+
	V_{\mathrm{NFW}}^2(R),
\end{equation}
where $V_{\mathrm{bar}}$ is computed from the SPARC baryonic decomposition
and $V_{\mathrm{NFW}}$ is the halo contribution derived from the catalogue
parameters.
In the analysis below, we adopt the ratio
$V_{\mathrm{bar}}/V_{\mathrm{NFW}}$ as an internal diagnostic proxy for
baryonic dominance. While model-dependent by construction, it provides a
self-consistent measure of the relative importance of the halo component
within the specific NFW implementation being tested.
\subsection{Implementation validation}
\label{subsec:nfw_validation}

Before performing the perturbation analysis, we verified that the NFW
velocity module reproduces the catalogue convention consistently. This
validation included inspection of the tabulated parameter definitions,
numerical confirmation of the $r_{200}$--$r_{\mathrm{s}}$ relation under the
adopted $V_{200}$ and $c_{200}$ convention, and a single-galaxy comparison
check against a well-sampled reference case. In particular, the validation
run for NGC\,3198 confirmed that the baseline NFW reconstruction and the
perturbed curves behave consistently with the catalogue parameterisation.

The NFW analysis therefore uses the published Li et al. baseline as its
reference state and does not introduce an independent halo refitting stage.

\subsection{Bounded perturbation scheme}
\label{subsec:nfw_perturbation}

To probe inner-region scatter sensitivity, we apply bounded perturbations to
the baseline NFW parameters rather than performing unrestricted refits. For
each galaxy, the halo mass scale and concentration are perturbed by
$\pm10$ and $\pm20$ per cent around their baseline values. These ranges are chosen to represent plausible bounded variation around calibrated fits without entering the regime of full re-optimisation. These
perturbations are applied separately to each parameter and jointly to both
parameters, subject to the requirement that the perturbed values remain
physically admissible.

The purpose of this procedure is not to re-optimise the halo model, but to
measure how strongly the adopted inner-region scatter metric responds to
controlled finite changes in a standard dark-matter implementation. The
resulting sensitivity is summarised by the maximum absolute scatter change
across the bounded perturbation set,
\begin{equation}
	\max |\Delta \sigma|_{\mathrm{NFW}}.
\end{equation}

\subsection{Results}
\label{subsec:nfw_results}

All 80 galaxies in the original H2 fleet were matched successfully to the
Li et al. NFW catalogue. After applying the common inner-region quality cut
described in Section~\ref{sec:data}, 74 galaxies remain in the final NFW
comparison sample.

The NFW perturbation analysis shows that bounded halo variations
exhibit a measurable inner-region scatter response across the sample. Across the 74-galaxy
comparison set, the median value of
$\max |\Delta \sigma|_{\mathrm{NFW}}$ is 0.0548 dex. Only one galaxy falls
below the 0.01 dex level, while 41 of the 74 systems exceed 0.05 dex. The
response therefore remains measurably non-zero for the majority of the
sample under the adopted perturbation scheme (Figure~\ref{fig:nfw_sensitivity}).

The suppression strength is also strongly modulated by baryonic dominance.
Using the median inner-region ratio
$V_{\mathrm{bar}}/V_{\mathrm{NFW}}$ for each galaxy, we find a strong
anti-correlation between this ratio and
$\max |\Delta \sigma|_{\mathrm{NFW}}$, with Spearman
$\rho = -0.899$ and $p = 1.8 \times 10^{-27}$. In other words, the bounded
NFW perturbations have less effect on the inner-region scatter metric when
the baryonic contribution dominates more strongly, and more effect when the
halo contribution is more comparable.

This trend is quantified by regime-stratified medians: the
baryon-dominated subsample shows median
$\max |\Delta \sigma|_{\mathrm{NFW}} = 0.0430$ dex, the balanced
subsample shows 0.0633 dex, and the dark-matter-dominated subsample
shows 0.0711 dex.

The NFW case therefore establishes a non-neutral comparison baseline for
the broader three-way analysis presented below.
\begin{figure}
	\includegraphics[width=\columnwidth]{nfw_scatter_sensitivity.png}
	\caption{Inner-region scatter sensitivity in the NFW perturbation analysis for the 74-galaxy comparison sample. The upper panel shows the maximum absolute scatter response, $\max |\Delta \sigma|_{\mathrm{NFW}}$, as a function of the median inner-region ratio $V_{\mathrm{bar}}/V_{\mathrm{NFW}}$. The lower panel shows the distribution of $\max |\Delta \sigma|_{\mathrm{NFW}}$ across the sample. The strong anti-correlation indicates that bounded halo perturbations have reduced leverage on the inner-region scatter metric as baryonic dominance increases.}
	\label{fig:nfw_sensitivity}
\end{figure}

\section{RAR perturbation analysis}
\label{sec:rar}

\subsection{RAR implementation}
\label{subsec:rar_implementation}

For the modified-gravity comparison, we use the RAR catalogue reported by
\citet{Li2020}. In this implementation, the baseline model is defined by
the adopted radial acceleration relation together with the tabulated
best-fitting galaxy-specific quantities, including the stellar
mass-to-light ratios and the geometric parameters used in the catalogue.
The acceleration scale $g_{\dagger}$ is fixed in the baseline
implementation, consistent with the convention adopted by the source
table.

For each galaxy, the RAR contribution is evaluated on the native SPARC
rotation-curve radii described in Section~\ref{sec:data}, so that the
inner-region scatter diagnostic is computed on the same observed points
used in the NFW and H2 comparisons. The resulting model velocity is then
used directly in the harmonised scatter metric defined in
Section~\ref{sec:methods}.

\subsection{Implementation validation}
\label{subsec:rar_validation}

Before carrying out the perturbation analysis, we verified that the RAR
velocity module reproduces the adopted catalogue convention consistently.
This validation included direct inspection of the source table,
confirmation of the fixed-$g_{\dagger}$ implementation, and a
single-galaxy reconstruction test using a well-sampled reference case.
In particular, the validation run for NGC\,3198 confirmed that the
baseline RAR reconstruction and the perturbed curves behave consistently
within the adopted implementation.

The RAR analysis therefore uses the published Li et al. baseline as its
reference state and does not introduce an independent refitting stage.

\subsection{Bounded perturbation scheme}
\label{subsec:rar_perturbation}

To probe inner-region scatter sensitivity within the adopted RAR
implementation, we apply bounded perturbations around the published
baseline solution rather than performing unrestricted refits. The
perturbation set includes variations of $\pm10$ and $\pm20$ per cent in
$g_{\dagger}$ and in the stellar disc mass-to-light ratio, together
with combined perturbations of these quantities.

These variations should be interpreted strictly as bounded sensitivity
tests within the adopted RAR implementation. They are not intended to
represent an exhaustive survey of MOND theory space, nor a full
re-optimisation of the RAR model for each galaxy. Their role here is
operational: to measure how strongly the same inner-region scatter
metric responds when the adopted RAR implementation is subjected to
controlled finite perturbations around its catalogue baseline.
We note that the treatment of $g_{\dagger}$ as a perturbed 
quantity in this context should not be interpreted as a test of 
MOND as a theory; $g_{\dagger}$ is treated as a universal 
constant in the standard RAR framework. The perturbations 
applied here probe the arithmetic sensitivity of the adopted 
implementation, not the physical variability of the acceleration 
scale \citep[cf.][]{Li2021}.
As in the NFW analysis, the resulting response is summarised by the
maximum absolute scatter change across the bounded perturbation set,
\begin{equation}
	\max |\Delta \sigma|_{\mathrm{RAR}}.
\end{equation}

\subsection{Results}
\label{subsec:rar_results}

The RAR perturbation analysis was successfully evaluated for 74 galaxies
in the common comparison sample, with the same six systems excluded by
the inner-region quality criterion described in
Section~\ref{sec:data}. The comparison is therefore performed on the
same 74-galaxy set used in the NFW and H2 analyses.

Across this common sample, the adopted RAR implementation shows an
intermediate level of inner-region scatter sensitivity. The median value
of $\max |\Delta \sigma|_{\mathrm{RAR}}$ is 0.0321 dex, which is clearly
above the near-neutral H2 response but below the NFW baseline. In this
sense, the RAR case occupies an intermediate position in the ranked
three-way comparison.

Unlike the NFW case, the scatter sensitivity in the adopted RAR implementation shows no significant correlation with the inner-region baryonic dominance ($\rho = +0.0485$, $p = 0.681$, $N = 74$). This indicates that the RAR response is a uniform property of the implementation's mapping between baryonic and total acceleration, rather than being modulated by the local halo-to-baryon ratio. The lack of a significant Spearman trend suggests that the RAR sensitivity is regime-independent within the tested perturbation bounds.

The RAR perturbation analysis therefore provides an intermediate
comparison case for the broader three-way analysis presented below: more
sensitive than the H2 implementation under the harmonised inner-region
diagnostic, but less sensitive than the NFW baseline
(Figure~\ref{fig:rar_sensitivity}).

\begin{figure}
	\includegraphics[width=\columnwidth]{rar_sensitivity.png}
	\caption{Inner-region scatter sensitivity in the adopted RAR
		implementation for the 74-galaxy comparison sample.
		\textit{Top panel:} Maximum absolute scatter change
		$\max|\Delta\sigma|_{\mathrm{RAR}}$ (dex) as a function of the
		inner-region baryonic-dominance proxy
		$V_{\mathrm{bar}}/V_{\mathrm{RAR}}$, with points coloured by
		regime classification.
		\textit{Bottom panel:} Distribution of
		$\max|\Delta\sigma|_{\mathrm{RAR}}$ across the common sample, shown
		with NFW and H2 reference levels for comparison. The adopted RAR
		implementation shows an intermediate sensitivity level, with median
		$\max|\Delta\sigma|_{\mathrm{RAR}} = 0.0321$ dex, placing it between
		the near-neutral H2 response and the NFW baseline under the same
		operational inner-region diagnostic.}
	\label{fig:rar_sensitivity}
\end{figure}


\section{Fixed-parameter kernel implementation}
\label{sec:h2}

\subsection{Operational result}
\label{subsec:h2_results}

The fixed-parameter H2 kernel implementation shows a near-neutral
inner-region response in the harmonised metric. The median absolute
response is $|\Delta \sigma|_{\mathrm{H2}} = 0.000478$~dex, with a
maximum value of 0.005323~dex. These values remain well below the
corresponding response levels measured for the RAR and NFW
implementations.

The H2 inner-region scatter response was computed using the explicit
$V_{\mathrm{total,H2}}$ output for all 74 galaxies in the common
comparison set.

This near-neutral behaviour distinguishes the H2 case from both
comparison classes considered above. It should not be interpreted here
as evidence that the underlying implementation lacks baseline fit
content. Rather, the present diagnostic isolates how the fixed-parameter
implementation responds to the adopted bounded perturbation protocol.
The purpose of the H2 result here is to provide the third element of
the comparative framework under the same inner-region metric, radial
cut, and operational definitions used elsewhere in the paper.


\section{Comparative analysis}
\label{sec:comparative}

\subsection{Ranked sensitivity pattern}
\label{subsec:ranked_pattern}

The common comparison set contains 74 galaxies evaluated under the same inner-region definition and the same operational scatter metric. This permits a direct comparison of bounded inner-region scatter sensitivity across the three implementations considered here.

A clear ranked pattern emerges. The H2 fixed-parameter kernel implementation shows the smallest response, with median absolute response $|\Delta \sigma|_{\mathrm{H2}} = 0.000478$~dex. The adopted RAR implementation shows an intermediate response, with median $\max |\Delta \sigma|_{\mathrm{RAR}} = 0.0321$~dex. The NFW analysis shows the largest sensitivity, with median $\max |\Delta \sigma|_{\mathrm{NFW}} = 0.0548$~dex.

The resulting hierarchy of inner-region scatter sensitivity is clearly defined as:
\begin{equation}
	|\Delta\sigma|_{\mathrm{H2}} \ll |\Delta\sigma|_{\mathrm{RAR}} < |\Delta\sigma|_{\mathrm{NFW}}.
\end{equation}
Table~\ref{tab:summary} summarises the comparative statistics for all three implementations.
\begin{table}
	\caption{Summary of inner-region scatter sensitivity across the three
		tested implementations on the common 74-galaxy comparison sample.}
	\label{tab:summary}
	\centering
	\scriptsize
	\setlength{\tabcolsep}{3pt}
	\begin{tabular}{p{0.27\columnwidth}p{0.17\columnwidth}p{0.20\columnwidth}p{0.18\columnwidth}}
		\toprule
		Property & NFW & RAR & H2 \\
		\midrule
		$N_{\mathrm{galaxies}}$ & 74 & 74 & 74 \\
		Median $\max|\Delta\sigma|$ (dex) & 0.0548 & 0.0321 & 0.000478 \\
		Spearman $\rho$ & $-0.899$ & $+0.0485$ & --- \\
		Spearman $p$-value & $1.8\times10^{-27}$ & $0.681$ & --- \\
		Correlation proxy & $V_{\mathrm{bar}}/V_{\mathrm{NFW}}$ & internal baryonic proxy & --- \\
		Perturbation parameters & $V_{200}$, $c_{200}$ & $g_{\dagger}$, $\Upsilon_{\mathrm{disc}}$ & fixed-parameter kernel \\
		\bottomrule
	\end{tabular}
	
	\vspace{2pt}
	\parbox{\columnwidth}{\footnotesize \textit{Note.} All values were computed using the same inner-region definition ($R < 0.5\,R_{\max}$) and the same log$_{10}$-RMS residual metric (dex). The H2 correlation is not reported because the near-neutral response does not provide a meaningful dynamic range for rank-correlation analysis.}
\end{table}
The H2 response is therefore more than an order of magnitude below the RAR level and approximately two orders of magnitude below the NFW baseline when measured in the same $\log_{10}$-RMS residual metric and on the same 74-galaxy comparison set.

This ranking does not by itself identify a preferred physical theory. Rather, it shows that bounded perturbations applied within these three tested implementations do not produce identical inner-region behaviour. The H2 case remains near-neutral, the adopted RAR implementation is more responsive, and the NFW baseline is more responsive still. The comparative result is therefore one of differing structural sensitivity, not simple equivalence under a common perturbation test.

Figure~\ref{fig:three_way_comparison} summarises this ranked pattern graphically for the common comparison sample.

\subsection{Structural dependencies}
\label{subsec:structural_dependencies}

Although the three implementations differ in overall sensitivity amplitude, they also differ in how that sensitivity tracks the local balance between baryonic and non-baryonic or non-Newtonian contributions. The clearest example is provided by the NFW case, for which the inner-region scatter response is strongly anti-correlated with the median inner-region ratio $V_{\mathrm{bar}}/V_{\mathrm{NFW}}$. Galaxies with stronger baryonic dominance exhibit smaller $\max |\Delta \sigma|_{\mathrm{NFW}}$, whereas systems with a more comparable halo contribution exhibit larger sensitivity. This indicates that the leverage of bounded halo perturbations is not uniform across the sample, but depends strongly on the relative importance of the additional component within the adopted radial domain.

The adopted RAR implementation displays a weaker overall sensitivity than NFW and does not show the same degree of regime separation in the present diagnostic. Its bounded perturbation response remains systematically above the near-neutral H2 case, but well below the NFW baseline. Within the adopted RAR implementation, the bounded scatter response is more constrained than in the perturbed NFW case, although no comparably strong monotonic trend is recovered in the present diagnostic.

The H2 implementation remains distinct in that its bounded inner-region scatter response stays close to zero across the full 74-galaxy comparison set. This does not imply that H2 is free of all structural limitations, nor that the three implementations can be reduced to a single common mechanism. It does show, however, that the degree of inner-region scatter sensitivity differs substantially across the tested implementations even when the metric, radial cut, and comparison set are held fixed.

Taken together, these results support a comparative structural reading of the inner-region diagnostic. Bounded perturbations tend to have more limited leverage where the additional component is less important, but the strength of that limitation differs substantially between the tested implementations. The NFW case provides the clearest baryonic-dominance trend, the adopted RAR implementation occupies an intermediate position, and the H2 implementation remains near-neutral under the same operational test.

\begin{figure}
	\includegraphics[width=0.95\columnwidth]{three_way_comparison.png}
	\caption{Three-way comparison of inner-region scatter sensitivity for the 74-galaxy common sample.
		\textit{Panel (a):} Maximum absolute scatter change $\max|\Delta\sigma|$ (dex) as a function of an inner-region baryonic-dominance proxy. For visual comparison, the horizontal placement is shown using the median inner-region ratio $V_{\rm bar}/V_{\rm NFW}$, with NFW points shown as orange circles ($N=74$), the adopted RAR implementation as blue diamonds ($N=74$), and H2 as green squares ($N=74$). Dashed and dotted horizontal lines indicate implementation medians. The orange solid line is a linear trend guide for the NFW distribution. The NFW Spearman rank correlation is computed with the same $V_{\rm bar}/V_{\rm NFW}$ proxy ($\rho=-0.899$, $p=1.8\times10^{-27}$, $N=74$), whereas the adopted RAR implementation shows no statistically significant monotonic correlation when evaluated against its own internal baryonic-dominance proxy ($\rho=+0.0485$, $p=0.681$, $N=74$).
		\textit{Panel (b):} Box plots of the $\max|\Delta\sigma|$ distributions for each implementation, showing median (white line), interquartile range (box), and 1.5\,IQR whiskers. All boxes reflect the full 74-galaxy sample. The ranked sensitivity pattern H2\,$<$\,RAR\,$<$\,NFW is evident across all summary statistics.}
	\label{fig:three_way_comparison}
\end{figure}

\section{Discussion}
\label{sec:discussion}
\subsection{Structural interpretation}
\label{subsec:structural_interpretation}

The comparative analysis shows that bounded inner-region scatter
sensitivity is not identical across the tested implementations. On the
common 74-galaxy basis used throughout the comparison, the three
implementations exhibit clearly different response patterns, with the
NFW baseline showing the largest bounded sensitivity, the adopted RAR
implementation showing a weaker intermediate response, and the H2
implementation remaining near-neutral. For context, typical observational uncertainties in SPARC rotation-curve
measurements correspond to scatter contributions of order 0.05--0.10 dex.
The NFW bounded response is therefore of observationally relevant scale,
whereas the H2 response remains substantially smaller under the adopted
bounded perturbation protocol.

The NFW analysis provides the clearest evidence for a structural
dependence on baryonic dominance. Its strong anti-correlation between
$\max |\Delta \sigma|_{\mathrm{NFW}}$ and the median inner-region ratio
$V_{\mathrm{bar}}/V_{\mathrm{NFW}}$ shows that bounded perturbations
have progressively less leverage as baryonic dominance increases. In
this sense, the NFW case behaves consistently with the basic
suppression picture outlined in Section~\ref{sec:methods}: where the
additional component is locally less important, finite perturbations of
that component have more limited influence on the inner-region scatter
metric.

By contrast, the adopted RAR implementation does not show a
statistically significant monotonic correlation between
$\max |\Delta \sigma|_{\mathrm{RAR}}$ and the chosen baryonic-dominance
proxy. This suggests that, within the adopted RAR implementation, the
bounded scatter response is comparatively uniform across the sample,
rather than showing the same strong baryonic-dominance dependence seen
in the NFW baseline.

Taken together, these results support a comparative structural reading
of the diagnostic. Bounded perturbations tend to have reduced leverage
in the inner regions, but the strength and sample dependence of that
reduction differ substantially across the tested implementations. The
NFW case provides the clearest baryonic-dominance trend, the adopted
RAR implementation occupies an intermediate position, and the H2
implementation remains near-neutral under the same operational test.

\subsection{Role of parameter constraints}
\label{subsec:parameter_constraints}

This possibility cannot be dismissed entirely, and it should be treated
as a real limitation of the comparison. The H2 implementation is indeed
more tightly bounded than the NFW and RAR perturbation sets in a
practical sense, and that stronger freezing may contribute to its
extreme neutrality. The present analysis therefore does not justify a
claim that the three implementations differ only because of their
underlying physical principles.

At the same time, the NFW result shows that the comparative pattern is
not reducible to parameter counting alone. The observed variance within
the NFW class itself---where the number of perturbed quantities remains
constant---demonstrates that the inner-region scatter sensitivity
varies strongly across the sample and tracks the degree of baryonic
dominance. This indicates that local structural conditions matter in
addition to parameter freedom.

This point should be interpreted carefully. The present comparison does
not establish a model-independent unit of perturbation amplitude across
the three implementations, and the resulting hierarchy should not be
read as a universal ranking of physical rigidity. The more limited
result is that, under bounded finite variations defined within each
implementation and evaluated with the same inner-region scatter metric,
the tested models exhibit measurably different response scales. In that
sense, the comparison is diagnostic rather than theory-adjudicating.

The present result does not imply that a near-neutral scatter response
is intrinsically or theoretically superior. It shows only that the tested
implementations exhibit measurably different bounded sensitivities
under the adopted inner-region diagnostic.

\subsection{Interpretation framework and limits}
\label{subsec:interpretation_limits}

The present study should be interpreted as a comparative bounded
sensitivity analysis, not as a general test of all dark-matter or
modified-gravity frameworks. The NFW case is included as a standard
halo-based baseline that is widely used in galaxy rotation-curve
modelling, but it does not represent the full scope of $\Lambda$CDM
phenomenology at galaxy scales. Likewise, the RAR analysis is confined
to the adopted implementation and perturbation protocol used here, and
should not be taken as an exhaustive test of MOND theory space. The H2 case is similarly restricted to the fixed-parameter
kernel implementation carried into the present comparison.

These limits are important for interpreting the comparative result
properly. The present paper does not establish a hierarchy that extends
beyond the tested implementations, nor does it identify a preferred
fundamental theory. Its contribution is narrower: it shows that, under
a harmonised inner-region diagnostic, specific bounded implementations
can exhibit clearly different scatter-sensitivity behaviour, and that
those differences may depend both on local baryonic dominance and on
the way the additional component enters the model construction.

\section{Conclusions}
\label{sec:conclusions}

We have presented a comparative analysis of inner-region scatter sensitivity across three distinct galaxy rotation-curve implementations using a common sample of 74 galaxies from the SPARC database. By adopting the same inner-region definition ($R < 0.5\,R_{\max}$), the same operational scatter metric (log$_{10}$-RMS residuals in dex), and a common bounded-perturbation framework, we have placed the NFW, RAR, and H2 implementations on a directly comparable diagnostic basis.

A clear ranked pattern emerges from this comparison. The fixed-parameter H2 kernel implementation shows a near-neutral response, with median $|\Delta \sigma|_{\mathrm{H2}} = 0.000478$~dex. The adopted RAR implementation occupies an intermediate position, with median $\max |\Delta \sigma|_{\mathrm{RAR}} = 0.0321$~dex, while the NFW baseline shows the largest bounded sensitivity, with median $\max |\Delta \sigma|_{\mathrm{NFW}} = 0.0548$~dex. The resulting ordering,
\[
|\Delta \sigma|_{\mathrm{H2}} \ll |\Delta \sigma|_{\mathrm{RAR}} < |\Delta \sigma|_{\mathrm{NFW}},
\]
shows that the three tested implementations do not produce equivalent inner-region scatter responses under the same operational diagnostic.

The NFW case also shows that inner-region sensitivity is strongly modulated by local baryonic dominance. The anti-correlation ($\rho = -0.899$, $p = 1.8\times10^{-27}$, $N=74$) between $\max |\Delta \sigma|_{\mathrm{NFW}}$ and the median inner-region ratio $V_{\mathrm{bar}}/V_{\mathrm{NFW}}$ indicates that bounded halo perturbations have progressively less leverage where the baryonic contribution is more important. By contrast, the adopted RAR implementation shows no statistically significant monotonic correlation when evaluated against its own internal baryonic-dominance proxy ($\rho = +0.0485$, $p = 0.681$, $N = 74$). The three tested implementations therefore differ not only in overall response amplitude but also in their structural dependence within the adopted diagnostic.

These differences should not be taken to identify a preferred physical theory. The present result is narrower and comparative: under the same inner-region definition, metric, and bounded-perturbation test, the three implementations exhibit distinct levels of scatter sensitivity. In this sense, the comparison isolates a measurable structural difference in how each implementation responds to controlled perturbations within the inner rotation-curve domain.

These conclusions should be interpreted within the scope of the tested implementations only. The NFW baseline does not represent all $\Lambda$CDM halo phenomenology, the adopted RAR implementation does not exhaust MOND theory space, and the H2 case is restricted to the fixed-parameter kernel implementation analysed here. Within those limits, the present work supports the use of inner-region scatter sensitivity as a comparative diagnostic for galaxy rotation-curve models and provides a reproducible benchmark for future bounded-perturbation studies.


\section*{Acknowledgements}
The rotation-curve data used in this work are from the SPARC database
\citep{Lelli2016}, which is publicly available at
\url{http://astroweb.cwru.edu/SPARC/}. The NFW and RAR catalogue
parameters are from \citet{Li2020}. The author thanks the SPARC team
for maintaining a publicly accessible, well-documented data set that
makes reproducible comparative analyses of this kind possible. This
analysis made use of \textsc{NumPy} \citep{numpy2020},
\textsc{SciPy} \citep{scipy2020}, \textsc{Matplotlib}
\citep{matplotlib2007}, and \textsc{Pandas} \citep{pandas2010}.
AI-assisted code review and manuscript editing support were used in the
preparation of this work. All scientific design, data analysis,
interpretation, and conclusions are the sole responsibility of the
author. This research received no specific grant from any funding
agency in the public, commercial, or not-for-profit sectors.
ORCID: Vitor M. F. Figueiredo \url{https://orcid.org/0009-0004-7358-4622}
\section*{Data availability}
The data underlying this article are available from the SPARC database
\citep{Lelli2016} and from the galaxy-by-galaxy NFW and RAR catalogue
fits of \citet{Li2020}. The analysis scripts, figure-generation code,
and reproducibility materials associated with this work are available in
the GitHub repository at
\url{https://github.com/VitorFigueiredoResearch/H2-SPARC-Comparative-Galaxy-Dynamics}
and are archived via Zenodo at
\url{https://doi.org/10.5281/zenodo.18793232}
\citep{FigueiredoH2zenodo2026}.
	
	\bibliographystyle{mnras}
	\bibliography{references}
	

\end{document}
